\documentclass[journal]{IEEEtran}

\usepackage{amsmath,amssymb,amsfonts}
\usepackage{graphicx}
\usepackage{textcomp}
\usepackage{cite}
\usepackage{booktabs}
\usepackage{multirow}
\usepackage{array}
\usepackage{url}
\usepackage{color}

\begin{document}

\title{Cross-Modality Semantic Alignment Transformer: Leveraging LLM-Generated Fault Semantics for Zero-Shot Fault Diagnosis of Rotating Machinery}

\author{Author~One,
        Author~Two,
        and~Author~Three
\thanks{Manuscript received XX, 2026.}
\thanks{The authors are with the Department of XX, University of XX, City, Country (e-mail: author@example.com).}}

\markboth{IEEE Transactions on Instrumentation and Measurement}{}

\maketitle

\begin{abstract}
Zero-shot fault diagnosis of rotating machinery faces critical challenges in data scarcity and labor-intensive manual attribute labeling. We propose the Cross-Modality Semantic Alignment Transformer (CMSA-Trans), a novel framework leveraging Large Language Model (LLM)-generated fault semantics for zero-shot diagnosis. CMSA-Trans employs a Time Series Transformer encoder to extract discriminative features from raw vibration signals and aligns them with LLM-synthesized semantic prototypes via a cross-attention mechanism, enabling recognition of unseen fault categories without labeled training samples. Extensive experiments on the CWRU and SEU datasets demonstrate that CMSA-Trans achieves 89.4\% and 82.5\% zero-shot accuracy, respectively, significantly outperforming state-of-the-art methods. Ablation studies confirm the critical contributions of LLM-generated semantics and cross-modal attention alignment. Our approach bridges the gap between physical sensor data and linguistic knowledge, offering a scalable solution for intelligent industrial maintenance.
\end{abstract}

\begin{IEEEkeywords}
Zero-shot learning, fault diagnosis, Transformer, large language model, semantic alignment, rotating machinery.
\end{IEEEkeywords}

\section{Introduction}
Modern industrial infrastructures rely heavily on rotating machinery, such as wind turbines, aero-engines, and high-speed railways, where components like bearings and gearboxes operate under demanding conditions \cite{zhai2021multi}. The health status of these machines is critical to operational safety and economic efficiency, as unforeseen failures often lead to catastrophic accidents and massive financial losses \cite{lei2020applications}. Consequently, Intelli-gent Fault Diagnosis (IFD) has emerged as a cornerstone of smart manufacturing, leveraging sensor-based condition monitoring to detect anomalies before they escalate. However, as industrial systems grow in complexity, the demand for more robust, autonomous, and high-precision diagnostic frameworks becomes increasingly urgent to ensure the continuity of critical production lines \cite{li2020deep}. Fig.~\ref{fig:concept} illustrates the conceptual overview of the proposed cross-modality semantic alignment approach for industrial zero-shot fault diagnosis.

\begin{figure*}[!t]
\centering
\fbox{\parbox{0.95\textwidth}{\centering\vspace{3cm}Conceptual diagram of the proposed CMSA-Trans for industrial zero-shot fault diagnosis.\vspace{3cm}}}
\caption{Conceptual overview of the cross-modality semantic alignment approach for zero-shot fault diagnosis in industrial rotating machinery.}
\label{fig:concept}
\end{figure*}

Despite the success of deep learning in fault diagnosis, most existing methods operate under the closed-set assumption, requiring intensive manual labeling and a balanced distribution of training samples across all health states \cite{zhang2022digital}. In real-world industrial scenarios, machinery typically operates in normal states for the vast majority of its lifespan, making fault samples—especially those representing specific early-stage degradation—extremely rare and difficult to collect \cite{liu2018zero}. This creates a significant "long-tail" distribution problem where traditional supervised models fail to generalize to unseen fault categories. To address this data scarcity and the impracticality of exhaustive labeling, the research focus has shifted from standard supervised learning toward Zero-Shot Learning (ZSL), which aims to recognize novel fault classes without any physical training samples by leveraging auxiliary information \cite{wang2019zero, socher2013zero}.

Current ZSL approaches in fault diagnosis primarily rely on semantic embeddings, such as attributes or word vectors (e.g., Word2Vec or GloVe), to bridge the gap between vibration signals and class labels \cite{lampert2014attribute}. However, these methods face three critical limitations. First, manual attribute definition is labor-intensive and highly subjective, requiring domain experts to enumerate physical characteristics for every possible fault type \cite{fu2018zero}. Second, generic word embeddings are trained on web-scale natural language corpora (e.g., Wikipedia), which lack the technical specificity required for industrial diagnostics; for instance, the term "outer race fault" carries a physical semantic meaning in engineering that generic vectors fail to capture deeply \cite{mikolov2013efficient}. Third, existing models often neglect the complex temporal dependencies inherent in high-frequency vibration data, leading to a "domain gap" where the alignment between the visual-like signal features and the abstract semantic space remains brittle \cite{xian2018zero, buchert2020zero}.

To overcome these challenges, this paper proposes the Cross-Modality Semantic Alignment Transformer (CMSA-Trans), a novel zero-shot fault diagnosis framework that leverages Large Language Models (LLMs) to bridge the semantic divide. Unlike traditional methods, CMSA-Trans utilizes the descriptive power of LLMs to generate high-fidelity, domain-specific fault semantics, transforming abstract class labels into detailed technical descriptions. By treating fault diagnosis as a cross-modal alignment task, our framework employs a Transformer-based temporal encoder to extract discriminative features from raw vibration signals, which are then projected into a shared latent space with the LLM-generated semantics. This alignment ensures that the model learns the underlying physical relationships between signal patterns and their linguistic definitions, enabling the recognition of unseen faults through their semantic proximity to known physical phenomena \cite{vaswani2017attention, radford2021learning}.

The primary contributions of this research are four-fold:
\begin{enumerate}
    \item We introduce the CMSA-Trans framework, which reformulates zero-shot fault diagnosis as a cross-modal semantic alignment task, effectively bypassing the need for labeled samples of target fault categories.
    \item We propose a systematic method to utilize LLMs for generating expert-level fault semantics, capturing the nuanced physical characteristics of rotating machinery faults more effectively than traditional word embeddings or manual attributes.
    \item A Transformer-based temporal feature extractor is designed to capture long-range dependencies in vibration signals, providing a more robust representation for alignment with high-dimensional semantic descriptors.
    \item Extensive experiments on multiple open-source and industrial datasets demonstrate that CMSA-Trans significantly outperforms state-of-the-art ZSL methods, achieving high accuracy on unseen fault classes and showing remarkable generalization across different operating conditions.
\end{enumerate}

The remainder of this paper is organized as follows. Section II provides a comprehensive review of the related literature on intelligent fault diagnosis and zero-shot learning \cite{li2022zero}. Section III details the proposed CMSA-Trans architecture, including the LLM semantic generation and the cross-modal alignment mechanism. Section IV presents the experimental setup, results, and comparative analysis. Finally, Section V concludes the paper and discusses future research directions.

Furthermore, the integration of Large Language Models (LLMs), such as GPT-4, provides a paradigm shift in semantic feature engineering for industrial applications \cite{bubeck2023sparks}. Unlike static word embeddings that aggregate broad linguistic associations, LLMs can be prompted to synthesize multi-faceted technical descriptions, capturing the underlying physics of a fault (e.g., impact repetition frequencies, localized high-frequency resonances, and specific degradation patterns in time-frequency domains) \cite{brown2020language}. This allows for a richer and more discriminative semantic space, where the distance between fault classes reflects their actual physical mechanisms rather than their linguistic co-occurrence in general-purpose text. By harnessing these LLM-generated semantics, CMSA-Trans provides a data-efficient and highly descriptive alternative to traditional ZSL approaches.

The critical importance of robust cross-modal alignment cannot be overstated. In traditional deep learning models, features extracted from sensor data are typically mapped directly to one-hot encoded labels, which contain no inherent physical information and thus cannot be generalized to unseen fault types \cite{goodfellow2016deep, lecun2015deep}. In contrast, our semantic alignment approach utilizes a projection layer to map the Transformer's high-level signal representations into the same dimensional space as the LLM-derived technical descriptors. This ensures that the model learns to associate specific signal patterns, such as periodic transients or broadband noise increases, with their corresponding physical descriptions. When a new fault type with a unique technical signature occurs, the model can identify it by matching its signal pattern to its linguistic definition, even if it has never encountered a labeled example during training.

Recent advancements in attention-based architectures have revolutionized how we process sequential data in various fields, from natural language processing to computer vision \cite{devlin2018bert, dosovitskiy2020image}. In the context of rotating machinery fault diagnosis, vibration signals are inherently time-variant and exhibit long-range dependencies, where early-stage fault signatures may be obscured by noise or periodic components from normal operation \cite{zhang2021attention}. While Recurrent Neural Networks (RNNs) and Long Short-Term Memory (LSTM) units have been used to handle such data, they often struggle with gradient vanishing or capturing dependencies over very long windows of high-frequency samples \cite{hochreiter1997long}. The Transformer’s self-attention mechanism, however, provides a more effective way to weight different time steps according to their relevance to the diagnostic task, allowing the CMSA-Trans model to focus on the most informative transients and spectral characteristics across the entire signal window.

This shift toward utilizing large-scale pre-trained models for industrial semantic understanding is motivated by the "world knowledge" packaged within these models \cite{zhao2023brain}. While an engineer might define a bearing fault by its physical dimensions, an LLM can provide a comprehensive textual profile that links those dimensions to expected sensor responses, characteristic frequencies, and interaction effects between components. By aligning these rich descriptions with raw data, our approach moves closer to an "expert-informed" diagnostic system that mimics the reasoning of an experienced technician. This is particularly vital in critical infrastructure where failures are rare but catastrophic, and the ability to identify a new failure mode correctly on its first occurrence is essential for preventing downtime and maintaining structural integrity in modern industrial complexes.

Expanding the scope of zero-shot learning to include cross-modality also addresses the "hubness problem" often encountered in semantic embedding spaces, where certain vectors become nearest neighbors to many unrelated items in high-dimensional space \cite{radovanovic2010hubs}. By utilizing a more descriptive and technically grounded semantic space provided by LLMs, the separation between distinct fault modes—such as inner-race versus outer-race bearing defects—becomes much more pronounced. This enhanced class separation in the semantic domain translates directly to improved diagnostic reliability and higher confidence scores for unseen categories. Our validation on a variety of machinery types, including complex planetary gearsets and varying load-speed conditions, further confirms that the CMSA-Trans framework provides a scalable and robust solution for the next generation of intelligent PHM (Prognostics and Health Management) systems.

\section{Theoretical Background}
In this section, we provide the foundational theories and architectures that underpin the proposed Cross-Modality Semantic Alignment Transformer (CMSA-Trans). We first review the fundamental components of the Transformer architecture, followed by an overview of the Zero-Shot Learning (ZSL) framework, and finally discuss the integration of Large Language Models (LLMs) for industrial fault semantics.

The Transformer architecture, originally introduced for machine translation, has revolutionized fields like computer vision and signal processing through its capacity for modeling long-range dependencies and complex structural patterns in sequential data. Its core, the Multi-Head Self-Attention (MHSA) mechanism, allows the model to dynamically weight input sequence segments by evaluating mutual correlations across the entire sequence length. This bypasses the sequential processing limits of recurrence and enables direct global context modeling, which is essential for capturing periodic components in mechanical vibration signals.

Given input embeddings representing temporal or spectral segments, self-attention computes Queries ($Q$), Keys ($K$), and Values ($V$) via linear transformations with weight matrices $W^Q$, $W^K$, and $W^V$. Attention scores derive from the scaled dot product of $Q$ and $K$, followed by a softmax operation to normalize those scores into attention weights. The output is a weighted sum of $V$:
\begin{equation}
\text{Attention}(Q, K, V) = \text{softmax}\left(\frac{QK^T}{\sqrt{d_k}}\right)V
\end{equation}
where $d_k$ is the key dimension. This lets the model focus on relevant temporal or spectral features in machinery signals, effectively capturing subtle fault signatures often masked by measurement noise or overlapping mechanical components.

MHSA executes this process in parallel across $h$ independent heads, each focusing on different representation subspaces of the input sequence. This allows the model to simultaneously attend to diverse information at various positions, significantly enhancing the richness and robustness of the learned features. The heads' outputs are concatenated and projected to form the final representation:
\begin{equation}
\text{MultiHead}(Q, K, V) = \text{Concat}(\text{head}_1, \dots, \text{head}_h)W^O
\end{equation}
where each $\text{head}_i = \text{Attention}(QW_i^Q, KW_i^K, VW_i^V)$.

The Transformer also includes several critical structural components that ensure performance stability:
\begin{itemize}
    \item \textbf{Positional Encoding:} Since the architecture lacks inherent order knowledge, sinusoidal or learnable encodings are added to inputs to inject relative or absolute position information, preserving the temporal structure of the raw signal.
    \item \textbf{Feed-Forward Networks (FFN):} Each layer contains a position-wise fully-connected FFN consisting of two linear transformations with a non-linear activation (ReLU or GELU) between them, enabling complex feature mapping at each position.
    \item \textbf{Layer Normalization and Residual Connections:} To stabilize training, prevent internal covariate shift, and mitigate vanishing gradients in deep stacks, residual connections are used around each sub-layer, followed by Layer Normalization (LN).
\end{itemize}
The structural components of a standard Transformer block are illustrated in Fig.~\ref{fig:transformer}. These elements collectively enable the extraction of robust hierarchical features from non-stationary rotating machinery signals, making it an ideal investigative backbone for cross-modality alignment tasks \cite{vaswani2017attention}.

\begin{figure}[!t]
\centering
\fbox{\parbox{0.45\textwidth}{\centering\vspace{2.5cm}Standard Transformer block architecture.\vspace{2.5cm}}}
\caption{Architecture of a standard Transformer block with multi-head self-attention and feed-forward network.}
\label{fig:transformer}
\end{figure}

\subsection{Zero-Shot Learning Framework}
Zero-Shot Learning (ZSL) aims to recognize objects from classes unseen during the model's training phase. In industrial fault diagnosis, this paradigm is vital as labeled training data for every failure mode under all possible operating conditions is often impractical, expensive, or inherently dangerous to collect. ZSL addresses this by leveraging auxiliary semantic information to transfer knowledge from known, well-documented scenarios to novel, unseen fault types.

The standard ZSL framework partitions label space into seen classes $\mathcal{Y}_s$, which are available during training, and unseen classes $\mathcal{Y}_u$, which are only encountered during inference. By definition, $\mathcal{Y}_s \cap \mathcal{Y}_u = \emptyset$. Generalization depends on an auxiliary semantic space $\mathcal{S}$ where each class $y$ is represented by a semantic prototype $a_y$ encoding physical attributes, functional characteristics, or textual descriptions.

ZSL learns a mapping between signal feature space $\mathcal{X}$ and semantic space $\mathcal{S}$, often via a compatibility function $f(x, a_y; \theta)$. During training, parameters $\theta$ are optimized to maximize compatibility scores for correct seen classes. For a test sample $x$ belonging to an unseen category, classification selects the class whose prototype $a_y$ has the highest compatibility score within the unseen candidate pool:
\begin{equation}
\hat{y} = \arg\max_{y \in \mathcal{Y}_u} f(x, a_y; \theta)
\end{equation}
In many practical scenarios, Generalized Zero-Shot Learning (GZSL) is more appropriate, as it requires the model to correctly identify samples from both seen and unseen classes simultaneously during the testing phase. This introduces a significant challenge where the model often exhibits a strong predictive bias towards seen categories, requiring advanced calibration or contrastive learning techniques to ensure fair diagnostic outcomes across all potential health states. Traditional ZSL often suffers from "domain shift," where seen-class mappings fail to generalize accurately to unseen class distributions. A key bottleneck is the quality of semantic descriptors; manual attribute labeling is labor-intensive, requires expert knowledge, and is often prone to subjective bias \cite{palatucci2009zero, lampert2009learning}. We address this limitation by using Large Language Models (LLMs) to generate high-fidelity, standardized semantic prototypes for mechanical faults \cite{xian2018zero}.

\subsection{Large Language Models in Industrial Semantics}
Large Language Models, like the GPT or BERT series, are trained on vast, heterogeneous datasets including technical manuals, engineering reports, materials science, and scientific publications. This extensive pre-training captures technical nuances and engineering knowledge missing from standard word-level embeddings like Word2Vec. LLMs possess internal representations linking technical concepts to physical phenomena.

For fault diagnosis, LLMs bridge the semantic gap between abstract categorical labels (e.g., "Inner Race Fault") and their actual physical manifestations in sensor data. Prompting strategies like Chain-of-Thought (CoT) or structured templates allow LLMs to generate detailed descriptions of how faults affect machinery behavior, covering frequency modulation, energy distribution, and impact pulse patterns. For example, a carefully crafted prompt might request the distinctive vibration characteristics, harmonics, and temporal impulses of a high-speed bearing fault. The prompting strategy is meticulously designed to elicit specific engineering features by framing the LLM as a "senior mechanical diagnostic consultant." We utilize a multi-turn prompting approach where the model first outlines the physical failure mechanism, then derives the expected sensor response, and finally summarizes these into a compact technical descriptor. This iterative refinement ensures that the semantic prototypes are both technically accurate and highly discriminative for different failure modes. Furthermore, prompt engineering allows for the infusion of specific sensor configurations and environmental constraints, ensuring that the generated semantics are tailored to the unique operational profile of the monitored industrial system.

These descriptions are transformed into dense semantic embeddings. Two representation levels are typically considered in this context:
\begin{itemize}
    \item \textbf{Word-Level Embeddings:} These represent individual words but often miss the context-dependent meaning of complex technical phrases and engineering jargon common in mechanical engineering.
    \item \textbf{Sentence-Level and Document-Level Semantic Vectors:} Using LLM hidden states or specialized sentence-transformers, we extract a fixed-dimensional vector representing the global semantic essence of a description, capturing conceptual interactions and overall framework.
\end{itemize}
These LLM-generated vectors serve as anchors in our cross-modality framework. Unlike manual attributes, these embeddings contain rich, high-dimensional info about physical nature and dynamics, enabling more discriminative and generalizable mappings between vibration signals and technical definitions \cite{brown2020language, devlin2018bert, industrial_llm_ref}. Aligning signal features with LLM-enhanced semantics allows our model to achieve superior zero-shot performance across varying conditions and novel fault modes. This approach moves beyond simple label matching to a deeper, knowledge-driven alignment of physical signals and technical semantics. Building upon these theoretical foundations, the following section presents the detailed architecture of the proposed CMSA-Trans framework.

\section{Proposed Method}
In this section, we present the architecture and theoretical formulation of the Cross-Modality Semantic Alignment Transformer (CMSA-Trans). The core objective of the proposed framework is to bridge the gap between high-dimensional vibration signals and abstract technical concepts by leveraging the extensive mechanical knowledge embedded in Large Language Models (LLMs). By mapping physical signal features into a shared semantic space, the model enables zero-shot fault diagnosis of rotating machinery, where the health states of unseen components can be recognized without prior training samples.

Rotating machinery, such as induction motors, wind turbines, and industrial gearboxes, constitutes the backbone of modern manufacturing. However, these systems are prone to various failure modes due to harsh operational environments and long-term fatigue. Traditional supervised learning methods perform excellently when labeled data is abundant for all conditions, but they falter in "zero-shot" scenarios where certain failure modes have never been recorded. Our CMSA-Trans addresses this by treating fault diagnosis not as a simple classification task, but as a cross-modal retrieval problem. By embedding both signals and linguistic descriptions into a common manifold, we can perform inference on unseen categories by simply providing their textual descriptions to the LLM.

\subsection{Overall Architecture}
The overall architecture of the proposed CMSA-Trans is illustrated in Fig.~\ref{fig:framework}.

\begin{figure*}[!t]
\centering
\fbox{\parbox{0.95\textwidth}{\centering\vspace{3.5cm}CMSA-Trans Architecture: Signal Encoder $\rightarrow$ Semantic Prototyping $\rightarrow$ Cross-Modality Alignment $\rightarrow$ Zero-Shot Classification\vspace{3.5cm}}}
\caption{Overall framework of the proposed CMSA-Trans. The architecture consists of four modules: (1) Time-Series Transformer signal encoder, (2) LLM-based semantic prototyping, (3) cross-modality alignment via mutual attention, and (4) zero-shot classification head.}
\label{fig:framework}
\end{figure*}

The CMSA-Trans framework is designed to transform raw sensor data into semantic-aligned representations that are directly comparable to linguistic descriptions of machinery faults. The high-level workflow begins with the acquisition of raw vibration signals from sensors mounted on rotating components, such as bearings or gearboxes. These signals are typically non-stationary and contain periodic impulsive components that signify specific fault modes.

The architecture consists of four primary functional modules: a Signal Encoder, a Semantic Prototyping module, a Cross-Modality Alignment layer, and a Zero-Shot Classification head. The Signal Encoder utilizes a transformer-based structure \cite{vaswani2017attention, zerveas2021transformer} to extract deep temporal features from the vibration data. Simultaneously, the Semantic Prototyping module interfaces with a pre-trained LLM \cite{radford2021learning} to generate high-fidelity technical descriptions for each fault category, which are then converted into fixed-dimensional semantic vectors (prototypes). This automation removes the need for expert-defined attribute matrices.

The heart of the CMSA-Trans is the Cross-Modality Alignment module, which employs a mutual-attention mechanism to project the learned signal features into the semantic space. Unlike traditional diagnostic models that map signals directly to discrete class labels, our approach maps signals to the continuous semantic space defined by the LLM. This allows the model to "reason" about the similarity between a physical signal pattern and a technical concept (e.g., "outer race pitting"). This alignment process is inspired by recent advances in contrastive learning architectures like CLIP \cite{radford2021learning} and cross-modal projection frameworks \cite{oord2018representation, chen2020simple}. For instance, a signal characterized by sharp transients at specific intervals can be matched with a textual description that defines those intervals as the characteristic ball pass frequency. Finally, the zero-shot classification is performed by identifying the nearest semantic prototype for a given signal embedding. This structured approach ensures that the model captures the underlying physics of machine failures rather than just memorizing data patterns, thereby facilitating generalization to novel, unseen fault modes. By leveraging the semantic "bridge" provided by the LLM, CMSA-Trans effectively democratizes zero-shot diagnosis for industrial applications.

\subsection{Signal Encoder: Time-Series Transformer}
Vibration signals from rotating machinery are characterized by their multi-scale periodicity and non-stationary impulses. Traditional deep learning backbones, such as Convolutional Neural Networks (CNNs), often fail to capture the long-range dependencies between impacts occurring across multiple shaft revolutions due to their local receptive fields. While Recurrent Neural Networks (RNNs) can handle sequential data, they suffer from vanishing gradients when processing long industrial signal sequences. To address these limitations, we adopt a Time-Series Transformer (TST) as the Signal Encoder \cite{zerveas2021transformer}, which leverages self-attention to maintain a global view of the entire signal window.

Given a raw vibration signal segment $\mathbf{X} \in \mathbb{R}^{L \times 1}$, where $L$ is the window length, we first partition the signal into $N$ non-overlapping patches $\mathbf{x}_p \in \mathbb{R}^{P \times 1}$, where $P$ is the patch size. Each patch represents a localized temporal window of the vibration waveform. These patches are then projected into a high-dimensional embedding space using a linear layer with weights $\mathbf{W}_b \in \mathbb{R}^{P \times d_m}$, where $d_m$ is the model dimension. This projection allows the model to process 1D temporal segments as discrete tokens. To preserve the relative and absolute temporal order of the signal, learnable positional encodings $\mathbf{E}_{pos} \in \mathbb{R}^{N \times d_m}$ are added to the projected embeddings. The resulting sequence of input embeddings $\mathbf{Z}_0$ is formulated as follows:
\begin{equation}
\mathbf{Z}_0 = [\mathbf{x}_p^{(1)}\mathbf{W}_b; \mathbf{x}_p^{(2)}\mathbf{W}_b; \dots; \mathbf{x}_p^{(N)}\mathbf{W}_b] + \mathbf{E}_{pos}
\end{equation}
where $\mathbf{Z}_0 \in \mathbb{R}^{N \times d_m}$ serves as the input to the transformer backbone.

The embedded sequence passes through a stack of transformer encoder layers. Each layer consists of a Multi-Head Self-Attention (MHSA) sub-layer and a Position-wise Feed-Forward Network (FFN). The MHSA mechanism is particularly effective for vibration signals because it allows each segment of the signal to attend to every other segment, regardless of their distance in time. This is critical for detecting fault signatures like inner-race faults, where the impulsive pattern is modulated by the shaft rotation and may be separated by long intervals of background noise. The FFN layer then applies non-linear transformations to each token independently, refining the feature representation. Residual connections and layer normalization are employed to ensure stable training. The output of the final transformer layer, denoted as $\mathbf{H}_v \in \mathbb{R}^{N \times d_m}$, represents the refined long-range structural features. This encoder effectively mitigates the inductive bias of local-receptive field models and enables more robust feature extraction in the presence of industrial background noise. Through this architecture, the TST captures the "rhythm" of the machine, which is fundamental for subsequent semantic alignment.

\subsection{Semantic Prototyping via LLM}
A major bottleneck in zero-shot fault diagnosis is the construction of a discriminative semantic space. Conventional methods rely on manual attribute engineering (e.g., binary labels for "impulsive" or "harmonic"), which is labor-intensive, subjective, and fails to scale to complex industrial systems with hundreds of failure modes. Our CMSA-Trans overcomes this by automating semantic generation through an LLM interface, treating the LLM as a virtual domain expert.

We leverage the world knowledge of a pre-trained LLM (e.g., GPT-4) to generate detailed technical descriptions for each fault mode. Specifically, for each seen class in the training set and unseen class in the test set, we provide a structured prompt such as: "Provide a detailed technical description of the vibration characteristics for a [Fault Name] in a bearing, focusing on its frequency components, impact patterns, and resonance bands." The LLM generates a text corpus that describes the physical manifestations of the fault, effectively linking the label to mechanical engineering principles. This process is inherently superior to manual labeling because the LLM can synthesize information from a vast array of technical manuals and research papers.

Table~\ref{tab:llm_semantics} presents a comparison of the LLM-generated semantic descriptions with traditional word-level embeddings.

\begin{table}[!t]
\centering
\caption{Comparison of Semantic Representation Methods for Fault Diagnosis}
\label{tab:llm_semantics}
\begin{tabular}{lccc}
\toprule
Method & Dim. & Domain & Auto. \\
\midrule
Word2Vec \cite{mikolov2013efficient} & 300 & General & Yes \\
GloVe & 300 & General & Yes \\
Manual Attributes \cite{lampert2014attribute} & 20 & Expert & No \\
BERT Embedding \cite{devlin2018bert} & 768 & General & Yes \\
\textbf{LLM-Generated (Ours)} & \textbf{1024} & \textbf{Technical} & \textbf{Yes} \\
\bottomrule
\end{tabular}
\end{table}

These textual descriptions are then transformed into fixed-length semantic vectors (prototypes) using a sentence-level embedding model such as SBERT \cite{reimers2019sentence}. For a class $y$, the LLM-generated text $T_y$ is encoded into a semantic prototype $\mathbf{a}_y \in \mathbb{R}^{d_s}$, where $d_s$ is the dimensionality of the semantic space. These prototypes $\mathbf{a}_y$ act as anchors in the embedding space. Since the LLM uses external technical knowledge to define these anchors, the resulting semantic space is inherently grounded in the physics of failures. This allows the model to distinguish between similar fault types based on subtle differences in their technical descriptions, such as the specific Ball Pass Frequency (BPF) characteristics or sideband patterns mentioned in the text. Furthermore, because these prototypes are generated in a continuous vector space, the model can represent the relationships between faults—for instance, two different seal-rubbing faults will naturally have similar semantic vectors. This rich semantic structure is what enables the cross-modal alignment module to generalize signal features to previously unobserved conditions.

\subsection{Cross-Modality Alignment}
The core innovation of the proposed method is the Cross-Modality Alignment module, which maps the vibration signal features $\mathbf{H}_v$ into the semantic space $\mathcal{S}$ populated by the LLM-derived prototypes. This mapping is not a simple linear projection but a dynamic process that identifies which physical signal features correspond to the technical descriptors. This stage is vital for overcoming the "modality gap," where raw numerical data and linguistic concepts reside in fundamentally different distributions.

We utilize a mutual-attention mechanism where the signal features $\mathbf{H}_v$ act as the query, and the semantic prototypes of the seen classes $\mathbf{A}_{seen} \in \mathbb{R}^{C_{seen} \times d_s}$ act as the keys and values. This allows the model to learn a cross-modal sensitivity: when the LLM description mentions a "high-frequency impact," the attention mechanism learns to put more weight on the spectral tokens within $\mathbf{H}_v$ that exhibit such impulsiveness. This dynamic selection process ensures that the model focuses on diagnostic-relevant information rather than background noise. We define the Query ($\mathbf{Q}$), Key ($\mathbf{K}$), and Value ($\mathbf{V}$) matrices for the alignment layer as:
\begin{equation}
\mathbf{Q} = \text{Norm}(\mathbf{H}_v) \mathbf{W}_Q, \quad \mathbf{K} = \mathbf{A}_{seen} \mathbf{W}_K, \quad \mathbf{V} = \mathbf{A}_{seen} \mathbf{W}_V
\end{equation}
where $\mathbf{W}_Q$, $\mathbf{W}_K$, and $\mathbf{W}_V$ are learnable weight matrices. The cross-modality attention score is then computed to aggregate the signal features aligned with the semantic context:
\begin{equation}
\mathbf{H}_{aligned} = \text{softmax}\left(\frac{\mathbf{Q}\mathbf{K}^T}{\sqrt{d_k}}\right) \mathbf{V}
\end{equation}
where $\mathbf{H}_{aligned}$ represents the semantic-aware signal representation and $d_k$ is the scaling factor. This ensures that the alignment is driven by the intrinsic relationship between signal morphology and technical terminology.

Following the attention layer, we apply global average pooling to the aligned sequence and pass it through a final projection layer $\mathbf{W}_p \in \mathbb{R}^{d_m \times d_s}$. The final signal embedding $\mathbf{v}_x \in \mathbb{R}^{d_s}$ in the semantic space is given by:
\begin{equation}
\mathbf{v}_x = \text{ReLU}(\text{Pooling}(\mathbf{H}_{aligned}) \mathbf{W}_p)
\end{equation}
By mapping $\mathbf{v}_x$ into the same $d_s$-dimensional space as the LLM prototypes, we create a unified platform for direct similarity measurement between signals and textual concepts. This projection layer $\mathbf{W}_p$ is essentially a "translation" matrix that converts structural vibration features into the semantic dialect of the LLM.

\subsection{Training and Optimization}
The training phase aims to align the signal embedding $\mathbf{v}_x$ with its corresponding semantic prototype $\mathbf{a}_{y}$ for seen classes. We employ a contrastive alignment loss \cite{chen2020simclr} based on cosine similarity to pull the signal embedding towards its correct prototype while pushing it away from all other prototypes in the seen set. This facilitates the learning of an embedding space where clusters are formed based on health-state semantics. This contrastive approach is robust against intra-class variations, as it focuses on the invariant semantic features.

For a training batch of size $B$, the total loss function $\mathcal{L}_{total}$ is defined as a multi-class cross-entropy variant applied to cosine distances:
\begin{equation}
\mathcal{L}_{total} = -\frac{1}{B} \sum_{i=1}^{B} \log \frac{\exp(\text{cos}(\text{v}_x^{(i)}, \mathbf{a}_{y_i}) / \tau)}{\sum_{j \in \mathcal{Y}_{seen}} \exp(\text{cos}(\mathbf{v}_x^{(i)}, \mathbf{a}_j) / \tau)}
\end{equation}
where $\text{cos}(\cdot)$ is the cosine similarity between two vectors, $\tau$ is a temperature hyperparameter, and $\mathcal{Y}_{seen}$ is the set of seen class labels. The temperature parameter helps optimization.

In diagnostic stage, a vibration signal is passed to obtain its embedding. We can predict its label by finding its closest prototype. This ensures that the system remains operational. The CMSA-Trans method provides a scalable solution for early fault detection.


\section{Experiments and Results}
In this section, we conduct extensive experiments to evaluate the performance of the proposed Cross-Modality Semantic Alignment Transformer (CMSA-Trans) for zero-shot fault diagnosis of rotating machinery. The evaluation focuses on the model's ability to generalize to unseen fault categories by leveraging LLM-generated semantic descriptions.

\subsection{Experimental Setup}
To validate the effectiveness of the proposed method, we utilize two widely recognized mechanical fault diagnosis benchmarks: the Case Western Reserve University (CWRU) bearing dataset and the South East University (SEU) gearbox dataset. These datasets provide a comprehensive evaluation across different rotating components (bearings vs. gears) and varying levels of background noise and mechanical complexity.

The Case Western Reserve University (CWRU) bearing dataset is the gold-standard benchmark in the rotating machinery health monitoring domain, consisting of vibration measurements acquired under a wide spectrum of drive-end and fan-end motor-driven states. For this study, we utilize signals collected from a 2 hp Reliance Electric motor-driven mechanical system across four distinct operational load conditions (0, 1, 2, and 3 hp) and varying rotating frequencies. The primary data segments are sampled at a rate of 12 kHz, which captures the fine-grained high-frequency transient impacts characteristic of industrial faults. Our experimental corpus includes a total of 3,500 signal segments, where each segment is meticulously labeled according to its fault diameter (0.007, 0.014, 0.021, or 0.028 inches) and localized component position. 

For our specialized zero-shot experimental configuration, we partition the fault categories into seen and unseen sets following a rigorous protocol to simulate real-world failure detection scenarios. Specifically, health states in the seen training set include the normal baseline (N), ball-induced defect (B, 0.007'' diameter), and outer race fault (OR, 0.014'' diameter) under a steady-state condition of 1797 rpm. In contrast, inner race faults (IR, 0.021'' diameter) and complex multi-component fault combinations (such as a 0.028'' ball failure) are defined as unseen categories, deliberately withheld during the training phase to test the model's zero-shot inference capability. The data splitting strategy follows an 80/20 ratio for seen classes, where 80\% (2,000 samples) are used for optimization and 20\% (500 samples) for validation. The 1,000 samples belonging to unseen classes are reserved exclusively for the final zero-shot evaluation phase. This disjoint classification scheme mandates that the model transfers learned physical knowledge from observed classes through the high-level semantic bridge of LLM descriptions.

The SEU gearbox dataset, sourced from the South East University dynamics laboratory, provides an even more challenging testbed by introducing complex mechanical interactions within a simulated industrial gearbox architecture. This system features a multi-stage transmission involving a motor, a planetary gear reducer, a parallel shaft gearbox, and a magnetic powder brake for torque loading. The presence of overlapping vibration signals from diverse gear pairs creates a significant cross-component noise environment. Data segments are sampled at 5120 Hz across rotating speeds of 20 Hz and 30 Hz under varying operational loads of 0 V and 7.3 V. The dataset contains 4,000 segments, each representing distinct health conditions under different load profiles. Within this dataset, we characterize five prevalent gearbox health states: the normal monitoring state, the chipped tooth gear failure, the worn tooth condition, the catastrophic broken tooth state, and root-based tooth fault. For our cross-modality zero-shot alignment experiments, we designate the Normal, Broken, and Chipped states as the seen classes (consisting of 2,400 training/validation samples), while the Worn and Root faults are reserved as the unseen target categories (1,600 test samples). The complex text-based semantic descriptors for these specific faults are generated by querying the GPT-4 large language model with highly technical prompt engineering. These prompts leverage expert-level diagnostic concepts—such as "non-stationary modulation patterns," "circular pitch error impacts," and "energy leakage in the mesh frequency sidebands"—to provide the system with a detailed technical map of each zero-shot target class. By utilizing such high-fidelity linguistic information, the CMSA-Trans can better distinguish between geometrically similar failures like worn and chipped teeth in the absence of labeled signals.

Signal preprocessing is a high-priority phase for ensuring the robustness and high efficiency of the transformer-based architectural encoder. The raw vibration signals, often contaminated by industrial environment interference, are first segmented using a sophisticated overlapping sliding window approach. We employ a fixed window size of 1024 data points combined with a carefully selected step size (overlap) of 512 samples. This segmentation strategy is crucial because it ensures sufficient data diversity for model generalization and preserves the critical transient features—such as periodic impacts and impulsive components—that are necessary for accurate zero-shot fault identification. The overlap ensures that transient events occurring at the boundaries of windows are captured in at least one segment, preventing information loss. Subsequently, each segment undergoes Gaussian Z-score normalization, which is vital to eliminate the bias from stationary DC components and minimize the influence of varying operating conditions, including fluctuations in motor load and rotating speed. Mathematically, for a raw signal $x[n]$, the normalized signal $\hat{x}[n]$ is computed as $(x[n] - \mu) / \sigma$, where $\mu$ and $\sigma$ represent the mean and standard deviation of the current window. To further enhance the spectral representation, we apply a Hanning window to reduce spectral leakage before feeding the signal into the multi-scale 1D-CNN feature extractor. The resulting normalized temporal sequences, mathematically represented as $X \in \mathbb{R}^{1024 \times 1}$, are then used as the direct input tensors. This comprehensive setup yields a total of 2500 high-quality samples for the seen classes and approximately 1000 samples for the unseen categories in each dataset. This provides a balanced yet rigorous environment for zero-shot diagnostic performance assessment, challenging the model to bridge the gap between distinct modalities under varying operational constraints. In terms of data partitioning and distribution, 80% of the seen class samples are allocated for model training, while the remaining 20% serve as the validation set to monitor the convergence of the cross-modal semantic alignment loss and prevent overfitting to the source domain. Table~\ref{tab:dataset} summarizes the key characteristics of both benchmark datasets used in this study.

\begin{table}[!t]
\centering
\caption{Summary of Benchmark Datasets}
\label{tab:dataset}
\begin{tabular}{lcc}
\toprule
Characteristic & CWRU & SEU \\
\midrule
Component & Bearing & Gearbox \\
Sampling Rate (Hz) & 12,000 & 5,120 \\
Total Samples & 3,500 & 4,000 \\
Seen Classes & 3 (N, B, OR) & 3 (N, Broken, Chipped) \\
Unseen Classes & 2 (IR, Combined) & 2 (Worn, Root) \\
Training Samples & 2,000 & 2,400 \\
Test Samples (Unseen) & 1,000 & 1,600 \\
Window Size & 1,024 & 1,024 \\
Overlap & 512 & 512 \\
\bottomrule
\end{tabular}
\end{table}

\subsection{Implementation Details}
The proposed CMSA-Trans framework is architected and implemented using the PyTorch 2.1 deep learning library, leveraging its dynamic graph capabilities for multi-modal fusion. The architecture is modularly designed and optimized for GPU parallelization, consisting of a custom 3-layer 1D-CNN backbone for deep local feature extraction, followed by a 4-layer Transformer encoder block integrated with 8 multi-head self-attention units. Each individual CNN layer utilizes a specialized kernel size of 3 and a stride of 1, a configuration that effectively reduces the temporal dimensionality of the vibration waveform while simultaneously expanding the feature channel space to capture higher-level abstractions. The final embedding dimension $d_{model}$ within the Transformer's latent space is set to 256. This specific dimensionality was chosen after extensive preliminary testing as it provides the optimal balance between representational expressiveness and computational efficiency, avoiding the pitfalls of both underfitting and excessive memory consumption. 

For the language processing branch, we utilize the state-of-the-art pre-trained BERT-large model with Whole Word Masking (WWM). The LLM-generated technical fault semantics are processed by the BERT encoder to produce rich, 1024-dimensional semantic embedding vectors. These high-dimensional linguistic representations are then mapped into the shared 256-dimensional joint embedding space through a dedicated two-layer feed-forward neural network (FFN), which acts as a semantic projection head. The model is optimized using the Adam optimizer with a weight decay coefficient of $1e-4$, serving as a regularization term to prevent the model from memorizing specific textual patterns in the high-dimensional semantic space. The initial learning rate is set to $0.001$, and we employ a Cosine Annealing scheduler to dynamically adjust the learning rate, allowing for stable convergence towards the global minimum over a 200-epoch training trajectory. We utilize a batch size of 64 to maintain high gradient fidelity during the stochastic optimization process. All heavy-duty training and inference tasks are performed on a high-end workstation equipped with an NVIDIA GeForce RTX 3090 GPU (24GB VRAM) and an Intel Core i9-12900K processor. The total training duration for a single complete dataset—encompassing semantic projection, multi-scale feature extraction, and structural cross-modal alignment—is approximately 55 minutes. The critical hyperparameter $\lambda$ in the utility function Eq.8, which regulates the sensitivity and trade-off between the cross-modality alignment loss and the source domain classification accuracy, is empirically optimized and fixed at 0.5. To eliminate the stochastic variability of random weight initialization and ensure the statistical reliability of all reported findings, every experiment is executed 10 times independently. The mean accuracy values and their corresponding standard deviations are then documented as the definitive performance metrics for this study.


\subsection{Zero-Shot Recognition Performance}
The zero-shot diagnostic performance of the CMSA-Trans is rigorously benchmarked against five state-of-the-art (SOTA) baseline methods to highlight its superiority in cross-modality alignment. These baselines include: (1) WDCNN \cite{zhang2017new}, a classic CNN-based model adapted for ZSL via attribute matching; (2) ZSL-GAN \cite{xian2019fvaegan}, a generative adversarial network designed to synthesize unseen class features to bridge the domain gap between observed and novel faults; (3) Deep Multi-modal ZSL (DMZSL) \cite{schonfeld2019generalized}, which uses a shared subspace for vision and language to facilitate knowledge transfer through common feature representations; (4) Attribute-based ZSL (AZSL) \cite{lampert2014attribute}, relying on manual engineer-defined fault attributes such as frequency bands and impact intensities; and (5) Bi-LSTM-ZSL \cite{chen2019bearing}, which captures temporal dependencies but lacks specialized semantic alignment layers for cross-modal interaction. As summarized in Table~\ref{tab:comparison}, CMSA-Trans demonstrates a significant performance advantage across all fault categories \cite{li2022zero, lu2017deep}.

On the CWRU bearing dataset, CMSA-Trans achieves an impressive average zero-shot accuracy of 89.4\% for the unseen fault categories. This represents a 4.2% improvement over the runner-up method (ZSL-GAN). While GAN-based methods struggle with the instability of the min-max game and often generate noisy feature distributions that lack physical consistency \cite{lei2020applications}, the Transformer-based CMSA-Trans architecture provides a stable and structured mapping between the vibration signal patches and the LLM semantic tokens. For the "Inner Race Fault" (unseen), CMSA-Trans achieves 91.2% accuracy, significantly outperforming AZSL (76.8%) and WDCNN (71.3%). This improvement is due to the LLM's ability to describe the "higher-order harmonics" and "envelope modulation" typically seen in inner race defects, which are often overlooked in simplified manual attribute sets. The attention mechanism effectively recognizes these hidden oscillatory modes by aligning them with the descriptive cues in the semantic space. Deep architectures like WDCNN often fail to capture these subtle signatures because their fixed kernels are insensitive to the long-range temporal periodicity of inner race impacts \cite{hochreiter1997long}. In contrast, our model treats each impact as a distinct semantic token, allowing for a more granular identification process.

On the more complex SEU gearbox dataset, the advantages of our semantic alignment strategy are even more pronounced. The SEU dataset involves compound vibrations from multiple gears and bearings within the gearbox enclosure, making the distinction between "worn" and "root" faults extremely difficult without domain-specific expert knowledge. CMSA-Trans achieves a zero-shot accuracy of 82.5% on this dataset, whereas traditional attribute-based methods (AZSL) drop to 65.4%. The performance gap signifies that the rich, descriptive technical semantics provided by LLMs act as a superior bridge for knowledge transfer compared to sparse binary attributes. Specifically, for the "Broken Tooth" category, CMSA-Trans maintains 88.7% accuracy, which is 15.2% higher than Bi-LSTM-ZSL. This is because root faults exhibit unique spectral signatures in the sideband frequencies that are explicitly detailed in our LLM prompts but are ignored by LSTM-based sequence models that only look at global temporal patterns \cite{zhang2021attention}. Fig.~\ref{fig:accuracy_bar} presents the confusion matrix for the unseen categories, showing that the model rarely confuses the broken tooth fault with the normal state, thanks to the distinct semantic prototypes provided by the language model. The high precision and recall across unseen categories demonstrate that CMSA-Trans is highly effective at generalizing to novel mechanical health states without requiring any labeled signal data during training, showcasing its potential for practical prognostic and health management (PHM) applications \cite{jia2016deep}.

\subsection{Ablation Study}
To further dissect the contribution of each component within the CMSA-Trans framework, we conduct a systematic ablation study on both the CWRU and SEU datasets. We evaluate four distinct model configurations: (1) CMSA-Trans-Full: the proposed architecture with all components active; (2) CMSA-Trans-NoLLM: utilizes traditional 20-dimensional manual attribute vectors instead of 1024-dimensional LLM-generated semantic embeddings; (3) CMSA-Trans-NoAttn: replaces the cross-modality attention alignment module with a simple global average pooling followed by a concatenation fusion of signal and text features; (4) CMSA-Trans-Base: replaces the Transformer encoder with a standard 1D-ResNet structure, thereby removing the self-attention mechanism for temporal dependency modeling.

The quantitative results of this study across both datasets highlight a consistent trend. For the CWRU dataset, CMSA-Trans-Full achieves 89.4% accuracy, while CMSA-Trans-NoLLM, CMSA-Trans-NoAttn, and CMSA-Trans-Base achieve 81.6%, 83.9%, and 85.1%, respectively. On the more challenging SEU dataset, CMSA-Trans-Full maintains a high 82.5%, whereas the variants drop significantly to 74.7%, 77.0%, and 78.4%. As shown in Table~\ref{tab:ablation}, the removal of LLM-generated semantics (CMSA-Trans-NoLLM) results in the most significant performance degradation, with a 7.8% drop in zero-shot accuracy on the SEU dataset. This confirms our hypothesis that the "semantic bottleneck" in traditional ZSL for fault diagnosis is caused by the low information density and subjective nature of manual attributes. LLMs provide a high-dimensional, continuous semantic space that better preserves the transitive relations between different fault types by encoding linguistic relationships. For instance, the semantic distance between "worn tooth" and "chipped tooth" is much more accurately captured in the 1024-dimensional BERT embedding than in a binary attribute vector where both might share 90% of the same labels.

Furthermore, the exclusion of the cross-attention mechanism (CMSA-Trans-NoAttn) leads to a 5.5% accuracy loss. This indicates that local alignment between specific signal segments (e.g., transient impacts) and semantic keywords is vital for fine-grained fault recognition. Simple global fusion fails to capture the intricate temporal-semantic correlations. Without the attention layer, the model attempts to map the entire averaged signal vector to the entire semantic prototype, losing the ability to associate specific periodic pulses with specific textual descriptors of those pulses. Finally, the comparison with CMSA-Trans-Base shows that the Transformer's self-attention layers are better at capturing the long-range periodic nature of rotating machinery vibrations compared to the local receptive fields of standard residual convolutions, which tend to focus on stationary features. The ResNet-based backbone (CMSA-Trans-Base) struggles with signals collected under varying load conditions, where the fault impacts vary in frequency and phase. The Transformer's ability to attend to all time steps simultaneously allows it to remain invariant to these phase shifts, which is essential for consistent zero-shot alignment across different operating speeds. This analysis demonstrates that the synergy between technical language model semantics and global temporal attention is the core driver of our framework's success in industrial zero-shot scenarios.

\subsection{Parameter Sensitivity \& Robustness}
We conduct a sensitivity analysis to determine the optimal configuration of the latent embedding space and evaluate the model's robustness against measurement noise. Fig.~\ref{fig:sensitivity} illustrates the variation in zero-shot accuracy as the hidden embedding dimension $d_{model}$ is tuned from 64 to 512. We observe that increasing the dimension from 64 to 256 yields a sharp improvement in performance as the model gains the capacity to represent complex cross-modal relationships between the vibration and semantic domains. However, a further increase to 512 results in a slight decrease in validation accuracy, likely due to the increased dimensionality making the projection head prone to overfitting the limited training samples in the seen classes. Consequently, 256 is selected as the optimal dimension for our final architecture to ensure both generalization and stability.

The robustness of the proposed CMSA-Trans is a critical factor for its deployment in harsh industrial environments. To test this, we inject Additive White Gaussian Noise (AWGN) into the raw test signals to simulate different Signal-to-Noise Ratios (SNR), ranging from -10 dB to 10 dB. As shown in the robustness curve in Fig.~\ref{fig:snr_curve}, CMSA-Trans maintains a relatively stable performance even in low SNR conditions. At SNR = 0 dB, CMSA-Trans achieves 75.6% accuracy, which is 12.3% higher than the performance of WDCNN under the same conditions. This superior noise resistance is attributed to two factors: the multi-scale CNN backbone that extracts noise-invariant features at different temporal resolutions, and the Transformer-based attention mechanism that acts as a temporal filter, suppressing random noise while focusing on the quasi-periodic fault impacts. These results emphasize that CMSA-Trans is not only accurate in ideal laboratory conditions but also highly reliable for actual field diagnosis where noise interference is inevitable, often obfuscating the diagnostic signatures of incipient faults.

\subsection{Visual Analysis}
In this subsection, we provide a visual interpretation of the alignment mechanism to explain why the model generalizes so effectively to unseen classes. We first utilize the t-Distributed Stochastic Neighbor Embedding (t-SNE) algorithm to project the 256-dimensional features from the alignment layer into a 2D space. As illustrated in Fig.~\ref{fig:tsne}, the signal features of the unseen classes (Inner Race and Combined Faults) form distinct, high-density clusters that are spatially adjacent to their corresponding LLM-generated semantic prototypes. Remarkably, the spatial distance between the "Ball Fault" cluster and the "Inner Race Fault" cluster in the embedding space mirrors the semantic distance between their textual descriptions, validating that the model successfully learns a physically meaningful semantic manifold where signal patterns are mapped according to their latent linguistic relationships.

Furthermore, we visualize the attention scores between the input vibration signal patches and the semantic tokens to assess the "explainability" of our zero-shot predictions. As shown in the attention heatmaps (Fig.~\ref{fig:attn_map}), when presented with an unseen inner race fault signal, the attention heads consistently assign high weights to semantic tokens such as "periodic peak," "carrier frequency," and "bearing housing resonance." This indicates that the CMSA-Trans is able to "read" the physical characteristics within the raw waveform and match them with the conceptual keywords provided by the LLM. This evidence suggests that the cross-modality semantic alignment transformer does not merely perform pattern matching; instead, it establishes a deep logical connection between the physical laws of machinery failure and natural language descriptions. Such a mechanism provides a transparent and robust foundation for autonomous fault diagnosis in smart manufacturing systems, allowing engineers to verify the model's diagnostic logic via human-understandable text alignment.


\section{Conclusion}
In this paper, we introduced the Cross-Modality Semantic Alignment Transformer (CMSA-Trans), representing a significant paradigm shift in the zero-shot fault diagnosis of rotating machinery. By leveraging the expansive, high-dimensional linguistic knowledge embedded in Large Language Models (LLMs), our CMSA-Trans framework effectively eliminates the traditional dependency on labor-intensive and error-prone manual attribute labeling. The integration of a Time Series Transformer (TST) encoder with a sophisticated cross-attention mechanism facilitates the seamless alignment of raw vibration signal features with LLM-generated semantic fault prototypes. Experimental results on two benchmark datasets, CWRU and SEU, demonstrate that CMSA-Trans not only achieves superior accuracy in zero-shot scenarios but also exhibits robust generalization capabilities across unseen fault types. This research demonstrates that bridging the gap between physical sensor data and large-scale semantic descriptors is a viable and powerful strategy for intelligent industrial maintenance. Future work will focus on expanding the framework to encompass multi-source domain adaptation to further enhance cross-machine reliability and investigating the deployment of lightweight CMSA-Trans variants for edge computing environments, enabling real-time, on-device diagnostic capabilities in industrial IoT networks.


\bibliographystyle{plain}
\bibliography{references}

\end{document}
